%% 전농사거리 소음측정계획서
%% 컴파일: xelatex (목차·전체 쪽수 반영을 위해 2회)
%% 본 문서는 생성형 AI(Claude)를 활용하여 작성되었음.
\documentclass[10.5pt,a4paper]{article}
\usepackage[a4paper,top=22mm,bottom=24mm,left=22mm,right=22mm,headsep=6mm,footskip=10mm]{geometry}
\usepackage{fontspec}
\usepackage{xetexko}
\IfFontExistsTF{Noto Serif CJK KR}{%
  \setmainfont{Noto Serif CJK KR}\setsansfont{Noto Sans CJK KR}%
  \setmainhangulfont{Noto Serif CJK KR}\setsanshangulfont{Noto Sans CJK KR}%
  \setmonohangulfont{Noto Sans CJK KR}[Scale=0.85]}{%
  \setmainfont{Times New Roman}\setsansfont{Apple SD Gothic Neo}%
  \setmainhangulfont{AppleMyungjo}[AutoFakeBold=2.5]\setsanshangulfont{Apple SD Gothic Neo}%
  \setmonohangulfont{Apple SD Gothic Neo}[Scale=0.85]}
\IfFontExistsTF{DejaVu Sans Mono}{\setmonofont{DejaVu Sans Mono}[Scale=0.82]}{%
  \setmonofont{DejaVuSansMono.ttf}[Scale=0.82,BoldFont=DejaVuSansMono-Bold.ttf]}

\usepackage{amsmath,amssymb}
\usepackage{array,booktabs,longtable,tabularx,makecell}
\usepackage[table]{xcolor}
\usepackage{graphicx}
\usepackage{enumitem}
\usepackage{needspace}
\usepackage{titlesec}
\usepackage{fancyhdr}
\usepackage{lastpage}
\usepackage[font=small,labelsep=colon,skip=4pt]{caption}
\usepackage[hidelinks]{hyperref}
\linespread{1.3}
\setlength{\parindent}{0pt}
\setlength{\parskip}{0.45em}
\setlength{\tabcolsep}{4pt}
\renewcommand{\arraystretch}{1.25}
\setlist{itemsep=0.15em,topsep=0.3em,leftmargin=1.6em}
\captionsetup[table]{name=표}
\captionsetup[figure]{name=그림}
\setcounter{secnumdepth}{3}
\setcounter{tocdepth}{2}
\definecolor{navy}{HTML}{1F3864}
\definecolor{headfill}{HTML}{D9E2F3}
\definecolor{inp}{HTML}{FFF8DC}
\definecolor{lab}{HTML}{F2F2F2}
\definecolor{team1}{HTML}{DDEBF7}
\definecolor{team2}{HTML}{E2EFDA}
\renewcommand{\thesection}{\arabic{section}.}
\renewcommand{\thesubsection}{\arabic{section}.\arabic{subsection}}
\renewcommand{\thesubsubsection}{\thesubsection.\arabic{subsubsection}}
\titleformat{\section}{\sffamily\Large\bfseries\color{navy}}{\thesection}{0.6em}{}
\titleformat{\subsection}{\sffamily\large\bfseries\color{navy}}{\thesubsection}{0.6em}{}
\titleformat{\subsubsection}{\sffamily\normalsize\bfseries}{\thesubsubsection}{0.6em}{}
\titlespacing*{\section}{0pt}{1.1em}{0.5em}
\titlespacing*{\subsection}{0pt}{0.9em}{0.4em}
\newcommand{\Laeq}{\ensuremath{L_{\mathrm{Aeq}}}}
\newcommand{\TBD}{\textcolor{gray}{[측정 후 기입]}}
\newcommand{\IN}{[\quad]}
\newcommand{\hd}[1]{\cellcolor{headfill}\textbf{\sffamily #1}}
\newcommand{\ip}{\cellcolor{inp}}
\newcommand{\lb}{\cellcolor{lab}}
\newcommand{\hr}[1]{\rule[-0.35\dimexpr#1\relax]{0pt}{#1}}
\newcommand{\notebox}[2][lab]{\par\noindent\colorbox{#1}{\parbox{\dimexpr\linewidth-2\fboxsep}{\small #2}}\par}
\renewcommand{\contentsname}{목차}
\pagestyle{fancy}\fancyhf{}
\fancyhead[L]{\small\sffamily 측정계획서}
\fancyhead[R]{\small\sffamily 전농사거리 소음현황조사}
\fancyfoot[L]{\scriptsize\sffamily\color{gray} 본 문서는 생성형 AI(Claude)를 활용하여 작성되었음}
\fancyfoot[R]{\small\sffamily \thepage\ / \pageref*{LastPage}}
\renewcommand{\headrulewidth}{0.4pt}
\fancypagestyle{plain}{\fancyhf{}\fancyfoot[L]{\scriptsize\sffamily\color{gray} 본 문서는 생성형 AI(Claude)를 활용하여 작성되었음}%
  \fancyfoot[R]{\small\sffamily \thepage\ / \pageref*{LastPage}}\renewcommand{\headrulewidth}{0pt}}
\begin{document}
\begin{titlepage}
\centering\vspace*{40mm}
{\sffamily\Huge\bfseries\color{navy} 전농사거리 소음현황조사 측정계획서\par}
\vspace{6mm}
{\sffamily\large 현재 소음도 조사와 소음지도 시뮬레이션 검증을 위한 현장 측정 계획\par}
\vspace{18mm}
{\renewcommand{\arraystretch}{1.5}
\begin{tabular}{>{\raggedright\arraybackslash}p{38mm}>{\raggedright\arraybackslash}p{90mm}}
\toprule
\cellcolor{headfill}\textbf{\sffamily 과목} & 소음진동설계 \\ \midrule
\cellcolor{headfill}\textbf{\sffamily 대상} & 전농사거리(사가정로 $\times$ 전농로) 주변 \\ \midrule
\cellcolor{headfill}\textbf{\sffamily 문서 상태} & 계획 v0.2 --- 측정점 2곳 확정(A안), 소음계 정보 미기입 \\ \midrule
\cellcolor{headfill}\textbf{\sffamily 작성일} & 2026-10-21 \\ \midrule
\cellcolor{headfill}\textbf{\sffamily 작성자} & \IN \\ \midrule
\cellcolor{headfill}\textbf{\sffamily 측정 예정일} & \IN\ (평일) \\
\bottomrule
\end{tabular}}
\vfill
\notebox{\textbf{본 문서는 생성형 AI(Claude)를 활용하여 작성되었음.} 기준값·조항은 국가법령정보센터 원문으로 확인한 것만 적었고, 원문에서 확인하지 못한 값은
`미확인'으로 표시하였다. 측정 결과는 포함하지 않는다.}
\end{titlepage}
\tableofcontents
\clearpage
\needspace{10\baselineskip}\section{조사 목적과 범위}
\needspace{8\baselineskip}\subsection{목적}
이번 소음현황조사는 두 가지 목적을 함께 만족하도록 계획한다. 같은 측정에서 두 목적의 자료를 얻되, 평균 방법과 비교 대상은 목적별로 나누어 처리한다.
\small
\begin{longtable}{>{\raggedright\arraybackslash}p{\dimexpr 0.2000\linewidth-2\tabcolsep\relax}>{\raggedright\arraybackslash}p{\dimexpr 0.3700\linewidth-2\tabcolsep\relax}>{\raggedright\arraybackslash}p{\dimexpr 0.4300\linewidth-2\tabcolsep\relax}}
\toprule
\hd{구분} & \hd{목적} & \hd{결과의 쓰임} \\ \midrule
\endfirsthead
\toprule
\hd{구분} & \hd{목적} & \hd{결과의 쓰임} \\ \midrule
\endhead
\bottomrule
\endlastfoot
가. 현재 소음도 조사 & 전농사거리 주변의 소음피해 현황을 국내 기준과 비교·평가 & 소음 환경기준(환경정책기본법 시행령 별표 1), 교통소음 관리기준(소음·진동관리법 시행규칙 별표 11) 초과 여부 \\ \midrule
나. 시뮬레이션 검증 & 측정값을 소음지도 소프트웨어(RLS-19 음원 + ISO 9613-2 전파) 계산값과 같은 위치·조건에서 비교 & 측정 \Laeq{}와 계산 \Laeq{}의 차이, 차이 원인(교통량·차종·반사·신호 대기 등) 검토 \\
\end{longtable}
\addtocounter{table}{-1}
\normalsize
\needspace{8\baselineskip}\subsection{범위}
\small
\begin{longtable}{>{\raggedright\arraybackslash}p{\dimexpr 0.1800\linewidth-2\tabcolsep\relax}>{\raggedright\arraybackslash}p{\dimexpr 0.8200\linewidth-2\tabcolsep\relax}}
\toprule
\hd{항목} & \hd{내용} \\ \midrule
\endfirsthead
\toprule
\hd{항목} & \hd{내용} \\ \midrule
\endhead
\bottomrule
\endlastfoot
공간 범위 & 교차로 중심(모델 설정값 EPSG:5186 X 205108.0, Y 553110.0) 주변. 측정점 2곳: A1 사가정로, A2 전농로(4장) \\ \midrule
인원·장비 & 2팀 $\times$ 3명(총 5--6명), 팀마다 소음계 1대. 두 측정점을 같은 시각에 측정(5장). 소음계 정보는 추후 기입 \\ \midrule
대상 소음 & 도로교통소음. 항공기·철도·건설작업 소음은 소음 환경기준 적용 대상이 아님(별표 1 비고 3) \\ \midrule
시간 범위 & 측정 예정일 \IN(평일, 월--금). 낮 4회(06:00--22:00), 밤 2회(22:00--06:00) --- 야간 측정 여부 \IN \\ \midrule
제외 & 실내 소음, 진동, 주파수 분석(장비가 지원하면 참고 기록만) \\
\end{longtable}
\addtocounter{table}{-1}
\normalsize
\needspace{8\baselineskip}\subsection{목적별 처리 구분}
\small
\begin{longtable}{>{\raggedright\arraybackslash}p{\dimexpr 0.1800\linewidth-2\tabcolsep\relax}>{\raggedright\arraybackslash}p{\dimexpr 0.4100\linewidth-2\tabcolsep\relax}>{\raggedright\arraybackslash}p{\dimexpr 0.4100\linewidth-2\tabcolsep\relax}}
\toprule
\hd{항목} & \hd{가. 법정 평가} & \hd{나. 시뮬레이션 비교} \\ \midrule
\endfirsthead
\toprule
\hd{항목} & \hd{가. 법정 평가} & \hd{나. 시뮬레이션 비교} \\ \midrule
\endhead
\bottomrule
\endlastfoot
측정방법 & ES 03301.1c, ES 03304.1d 그대로 & 같은 측정 조건 + 동시 교통량·속도·기상 기록 \\ \midrule
1회 측정시간 & 회차당 30분(10분·5분 이상 요건 충족) & 같은 구간 \\ \midrule
여러 회 평균 & \textbf{산술평균}(원문 규정) & \textbf{에너지평균}(본 조사 자체 설정, 2.4) \\ \midrule
배경소음 보정 & 하지 않음(적용 방법에 조항 없음) & 보정하지 않은 값과 이상소음 구간 표시를 함께 보고 \\ \midrule
비교 대상 & 별표 1·별표 11 기준값 & 같은 시간대·회차의 동시 교통량으로 계산한 모델 \Laeq \\ \midrule
최종 자릿수 & 정수(소수 첫째 자리에서 반올림) & 소수 첫째 자리 \\
\end{longtable}
\addtocounter{table}{-1}
\normalsize
\needspace{10\baselineskip}\section{적용 기준 및 측정방법}
아래 내용은 국가법령정보센터(law.go.kr)의 현행본 원문에서 확인한 것이다(2026-10-19·21). 공정시험기준 원문 PDF는 \texttt{field\_survey/references/}에 보관하였다. 원문에 규정이 없거나 원문을 보지 못한 항목은 `미확인'으로 적었다.
\needspace{8\baselineskip}\subsection{적용 문서와 현행본}
\small
\begin{longtable}{>{\raggedright\arraybackslash}p{\dimexpr 0.3600\linewidth-2\tabcolsep\relax}>{\raggedright\arraybackslash}p{\dimexpr 0.5000\linewidth-2\tabcolsep\relax}>{\raggedright\arraybackslash}p{\dimexpr 0.1400\linewidth-2\tabcolsep\relax}}
\toprule
\hd{문서} & \hd{현행본} & \hd{확인} \\ \midrule
\endfirsthead
\toprule
\hd{문서} & \hd{현행본} & \hd{확인} \\ \midrule
\endhead
\bottomrule
\endlastfoot
소음·진동공정시험기준 & 국립환경과학원고시 제2026-39호(2026-09-14 일부개정·시행), 총 17종 & 원문 \\ \midrule
ES 03300.c 총칙(소음) & 2026년판 & 원문 \\ \midrule
ES 03301.1c 환경기준 중 소음측정방법 & 2020년판 & 원문 \\ \midrule
ES 03304.1d 도로교통소음관리기준 측정방법 & 2024년판 & 원문 \\ \midrule
환경정책기본법 시행령 [별표 1] 환경기준 & 대통령령 제36726호(2026-10-08 시행), 별표 1은 2025-10-01 개정분 & 원문 \\ \midrule
소음·진동관리법 시행규칙 [별표 11] 교통소음·진동의 관리기준 & 기후에너지환경부령 제42호(2026-06-22 시행), 별표 11은 2019-12-31 개정분 & 원문 \\ \midrule
소음지도의 작성방법(측정점 선정·검증 기준) & 기후에너지환경부고시 제2026-19호(2026-01-18 일부개정·시행), 별표 2 & 원문 \\ \midrule
KS C IEC 61672-1(소음계 규격) & 공정시험기준 참고자료에 2022년판으로 인용. KS 본문은 열람하지 않음 & 2차 자료 \\ \midrule
ISO 1996-2 & 원문 열람하지 않음 & 미확인 \\
\end{longtable}
\addtocounter{table}{-1}
\normalsize
\needspace{8\baselineskip}\subsection{측정방법 요구사항}
\footnotesize
\begin{longtable}{>{\raggedright\arraybackslash}p{\dimexpr 0.1700\linewidth-2\tabcolsep\relax}>{\raggedright\arraybackslash}p{\dimexpr 0.4700\linewidth-2\tabcolsep\relax}>{\raggedright\arraybackslash}p{\dimexpr 0.2500\linewidth-2\tabcolsep\relax}>{\raggedright\arraybackslash}p{\dimexpr 0.1100\linewidth-2\tabcolsep\relax}}
\caption{측정방법 요구사항(원문 확인 결과)}\\
\toprule
\hd{항목} & \hd{요구사항 요약} & \hd{출처} & \hd{확인} \\ \midrule
\endfirsthead
\toprule
\hd{항목} & \hd{요구사항 요약} & \hd{출처} & \hd{확인} \\ \midrule
\endhead
\bottomrule
\endlastfoot
소음계 & KS C IEC 61672-1 클래스 2 이상 & ES 03301.1c 3.1, ES 03304.1d 3.1 & 원문 \\ \midrule
소음계 성능 & 주파수 31.5 Hz--8 kHz 이상, 35--130 dB 이상, 레벨레인지 전환오차·눈금오차 0.5 dB 이내 & ES 03300.c 3.3.3 & 원문 \\ \midrule
교정 & 전원·동작 점검과 매회 교정. 표준음 발생기 오차 $\pm$1 dB 이내 & ES 03301.1c 3.2.2, ES 03300.c 3.2.3 & 원문 \\ \midrule
교정 전후 허용 편차 & 적용 방법에 규정 없음 & --- & 미확인 \\ \midrule
청감보정·동특성 & A 특성, 빠름(fast) & 두 방법 3.3 & 원문 \\ \midrule
샘플주기 & 1초 이내(이하) & ES 03301.1c 7.1.1, ES 03304.1d 3.3.3 & 원문 \\ \midrule
측정점(도로변지역) & 건축물에서 도로 방향 1.0 m, 지면 위 1.2--1.5 m. 보도 없으면 도로단 & ES 03301.1c 5.1.3 & 원문 \\ \midrule
도로변지역 범위 & 도로단에서 차선수 $\times$ 10 m & ES 03301.1c 5.1.1 주 1 & 원문 \\ \midrule
측정점(도로교통) & 피해 예상 거주 건물 경계에서 1.0--3.5 m 중 소음도 높을 곳, 1.2--1.5 m. 2층 이상은 창문 밖 0.5--1.0 m 가능, 창문 경계면 측정 시 +1.5 dB(A) & ES 03304.1d 5.1 & 원문 \\ \midrule
마이크로폰 & 삼각대 원칙, 손으로 들면 몸에서 0.5 m 이상, 주 소음원 방향 & 두 방법 5.2.1 & 원문 \\ \midrule
기상 & 풍속 2 m/s 이상 방풍망 필수, 5 m/s 초과 측정 금지 & 두 방법 5.2.1.4 & 원문 \\ \midrule
강우·노면 제한 & 적용 방법에 규정 없음 & --- & 미확인 \\ \midrule
측정일 & 평일(월--금). 도로교통은 주말 통행 증가 시 주말 가능 & 두 방법 5.2.2 & 원문 \\ \midrule
횟수(환경기준) & 낮: 2시간 이상 간격 4회 이상. 밤: 2시간 간격 2회 이상. 1회 5분 이상 & ES 03301.1c 5.3, 7.1.1 & 원문 \\ \midrule
횟수(도로교통) & 지점 2개 이상. 주간·야간 각각 4시간 이상 간격 2회 이상(야간은 00--04시 포함). 1회 10분 이상 & ES 03304.1d 5.3, 7.1.1 & 원문 \\ \midrule
평균 & 산술평균. 도로교통은 지점별 산술평균 후 지점들을 산술평균 & ES 03301.1c 5.3, ES 03304.1d 7.1.1.2 & 원문 \\ \midrule
자릿수 & 계산 중 소수 첫째 자리, 최종값은 반올림해 정수 & 두 방법 7.1 & 원문 \\ \midrule
배경소음 보정 & 두 방법에 조항 없음(보정표는 생활소음 방법 ES 03303.1d에만 있음) & ES 03303.1d 7.2 & 원문(미적용) \\ \midrule
기록 서식 & [서식 1] 환경소음 측정자료 평가표, [서식 6] 도로교통소음 측정자료 평가표 & 두 방법 10.1 & 원문 \\
\end{longtable}
\normalsize
\needspace{8\baselineskip}\subsection{기준값}
\small
\begin{longtable}{>{\centering\arraybackslash}p{\dimexpr 0.2200\linewidth-2\tabcolsep\relax}>{\centering\arraybackslash}p{\dimexpr 0.3000\linewidth-2\tabcolsep\relax}>{\centering\arraybackslash}p{\dimexpr 0.2400\linewidth-2\tabcolsep\relax}>{\centering\arraybackslash}p{\dimexpr 0.2400\linewidth-2\tabcolsep\relax}}
\caption{소음 환경기준(환경정책기본법 시행령 별표 1, 단위 $L_{\mathrm{eq}}$ dB(A))}\\
\toprule
\hd{지역 구분} & \hd{적용 대상지역} & \hd{낮(06:00--22:00)} & \hd{밤(22:00--06:00)} \\ \midrule
\endfirsthead
\toprule
\hd{지역 구분} & \hd{적용 대상지역} & \hd{낮(06:00--22:00)} & \hd{밤(22:00--06:00)} \\ \midrule
\endhead
\bottomrule
\endlastfoot
일반지역 & ``가''지역 & 50 & 40 \\ \midrule
일반지역 & ``나''지역 & 55 & 45 \\ \midrule
일반지역 & ``다''지역 & 65 & 55 \\ \midrule
일반지역 & ``라''지역 & 70 & 65 \\ \midrule
도로변지역 & ``가'' 및 ``나''지역 & 65 & 55 \\ \midrule
도로변지역 & ``다''지역 & 70 & 60 \\ \midrule
도로변지역 & ``라''지역 & 75 & 70 \\
\end{longtable}
\normalsize
{\small ``가'' 녹지·보전관리·농림·자연환경보전·전용주거지역, 종합병원·학교·공공/특수도서관 부지경계 50 m 이내. ``나'' 생산관리·일반주거·준주거지역. ``다'' 상업·계획관리·준공업지역. ``라'' 전용공업·일반공업지역.\par}
\small
\begin{longtable}{>{\raggedright\arraybackslash}p{\dimexpr 0.6400\linewidth-2\tabcolsep\relax}>{\centering\arraybackslash}p{\dimexpr 0.1800\linewidth-2\tabcolsep\relax}>{\centering\arraybackslash}p{\dimexpr 0.1800\linewidth-2\tabcolsep\relax}}
\caption{교통소음 관리기준 --- 도로(시행규칙 별표 11 제1호, $L_{\mathrm{eq}}$ dB(A))}\\
\toprule
\hd{대상 지역} & \hd{주간(06:00--22:00)} & \hd{야간(22:00--06:00)} \\ \midrule
\endfirsthead
\toprule
\hd{대상 지역} & \hd{주간(06:00--22:00)} & \hd{야간(22:00--06:00)} \\ \midrule
\endhead
\bottomrule
\endlastfoot
가. 주거·녹지·보전관리지역, 관리지역 중 취락·주거개발진흥·관광휴양개발진흥지구, 자연환경보전지역, 학교·병원·공공도서관·입소 100명 이상 노인의료복지시설·영유아보육시설 부지경계 50 m 이내 & 68 & 58 \\ \midrule
나. 상업·공업·농림지역, 관리지역 중 산업유통개발진흥지구 및 가목 외 관리지역, 미고시지역 & 73 & 63 \\
\end{longtable}
\normalsize
\needspace{8\baselineskip}\subsection{평균 방법}
\textbf{법정 평가(목적 가)}는 원문 규정대로 같은 측정점의 여러 회차 측정값을 산술평균한다(ES 03301.1c 5.3, ES 03304.1d 7.1.1.2).
\begin{equation*} L_{P} = \frac{1}{n}\sum_{i=1}^{n} L_{\mathrm{Aeq},i} \end{equation*}
\textbf{시뮬레이션 비교(목적 나)}는 측정 구간 길이 $T_i$로 가중한 에너지평균을 쓴다. 등가소음도의 정의(ES 03300.c 2.10)에서 나온 식이며, 여러 회 측정값을 이렇게 합치는 것은 법정 방법이 아닌 본 조사의 자체 설정이다.
\begin{equation*} L_{\mathrm{Aeq},T} = 10\lg\!\left(\frac{\sum_i T_i\,10^{0.1L_{\mathrm{Aeq},i}}}{\sum_i T_i}\right) \end{equation*}
\notebox{두 평균은 값이 다르다. 측정값의 변동이 클수록 에너지평균이 산술평균보다 커진다. 보고서에는 두 값을 섞지 않고 목적별 표에 따로 적는다.}
\needspace{8\baselineskip}\subsection{소음지도 검증 기준(「소음지도의 작성방법」 별표 2)}
\footnotesize
\begin{longtable}{>{\raggedright\arraybackslash}p{\dimexpr 0.2000\linewidth-2\tabcolsep\relax}>{\raggedright\arraybackslash}p{\dimexpr 0.4700\linewidth-2\tabcolsep\relax}>{\raggedright\arraybackslash}p{\dimexpr 0.3300\linewidth-2\tabcolsep\relax}}
\toprule
\hd{항목} & \hd{고시 내용} & \hd{이번 조사 적용} \\ \midrule
\endfirsthead
\toprule
\hd{항목} & \hd{고시 내용} & \hd{이번 조사 적용} \\ \midrule
\endhead
\bottomrule
\endlastfoot
검증 기준(제15호) & 계산값과 실측값의 오차 평균 $\pm$3 dB, 표준편차 3을 넘지 않는 것을 원칙. 지점별 오차 $\pm$5 dB 이상이면 원인 분석 & 보고서 6장 판정 기준 \\ \midrule
실측 지점(제15호) & 1시간 교통량 1,000대 이상인 도로당 한 지점 이상. 교통소음한도 초과 건물·공동주택 단지의 50\,\% 이상 지점 & 사가정로·전농로 각 1지점(A1·A2). 50\,\% 요건은 2지점으로 충족 못 함 → 한계로 보고 \\ \midrule
외벽 예측 위치(제14호 라목) & 건물에서 소음원 방향 1.0 m, 건물 반사음 포함 & 측정점을 외벽 앞 1.0 m에 두고 모델 점 수음점도 같은 위치 \\ \midrule
평가 단위(제10호) & 주간(06--22시)·야간(22--06시), 등가소음도 $L_{\mathrm{eq}}$ dB(A) & 같음 \\ \midrule
차종 구분(제12호 다목) & RLS: 소형(승용·소형트럭·중형트럭 3.5 t 미만), 대형(중형·대형버스, 대형트럭 3.5 t 이상). 트레일러 없으면 대형1, 있으면 대형2 & 계수 차종 대응(7.1) \\
\end{longtable}
\addtocounter{table}{-1}
\normalsize
\notebox[inp]{교통량 1,000대/h 조건: 교통량 산정근거 문서 값으로 낮 시간 사가정로 약 1,260대/h, 전농로 약 1,580대/h로 추정되어 두 도로 모두 해당한다(추정값, 실측 아님). 2지점이면 비교 수가 적으므로, 모델을 회차마다 그 회차의 동시 교통량으로 계산해 측정점 $\times$ 회차(최대 2 $\times$ 6 = 12개)로도 비교한다(자체 설정).}
\needspace{8\baselineskip}\subsection{미확인 항목}
\begin{itemize}
\item 측정점 A1·A2 필지의 용도지역 → 적용 기준값(가--라, 별표 11 가·나) 미정(3.3)
\item 교정 전후 허용 편차, 강우·노면 조건: 적용 방법 원문에 규정 없음. 자체 운영 기준은 6장에 `자체'로 표시
\item KS C IEC 61672-1 본문, ISO 1996-2 본문: 열람하지 않음
\end{itemize}
\needspace{10\baselineskip}\section{조사 대상 지역}
\needspace{8\baselineskip}\subsection{위치}
전농사거리는 서울특별시 동대문구의 사가정로와 전농로가 만나는 교차로이다. 소음지도 모델의 계산 중심은 EPSG:5186 좌표 X 205108.0, Y 553110.0(모델 설정 파일 기준)이며, 경위도는 약 북위 37.5776°, 동경 127.0580°(수치지도 객체 좌표에서 읽은 값)이다.
\needspace{8\baselineskip}\subsection{주변 도로}
\small
\begin{longtable}{>{\raggedright\arraybackslash}p{\dimexpr 0.2000\linewidth-2\tabcolsep\relax}>{\centering\arraybackslash}p{\dimexpr 0.1800\linewidth-2\tabcolsep\relax}>{\centering\arraybackslash}p{\dimexpr 0.1600\linewidth-2\tabcolsep\relax}>{\centering\arraybackslash}p{\dimexpr 0.2200\linewidth-2\tabcolsep\relax}>{\raggedright\arraybackslash}p{\dimexpr 0.2400\linewidth-2\tabcolsep\relax}}
\toprule
\hd{도로명} & \hd{차로수(수치지도)} & \hd{도로폭(m)} & \hd{교차로 중심과 거리} & \hd{모델 교통량} \\ \midrule
\endfirsthead
\toprule
\hd{도로명} & \hd{차로수(수치지도)} & \hd{도로폭(m)} & \hd{교차로 중심과 거리} & \hd{모델 교통량} \\ \midrule
\endhead
\bottomrule
\endlastfoot
사가정로 & 4--6 & 24.8--31.0 & 교차로 통과 & 가정값(실측 아님) \\ \midrule
전농로 & 2--5 & 19.8--25.5 & 교차로 통과 & 가정값(실측 아님) \\ \midrule
전농로25길 & 5 & 25.1--25.8 & 교차로 북측 접근로 & 가정값(표준값) \\ \midrule
사가정로5길 & 1(일방)·2 & 5.1--17.1 & 약 170 m-- & 가정값(표준값) \\ \midrule
사가정로17길 & 2 & 15.0 & 약 216 m-- & 가정값(표준값) \\
\end{longtable}
\addtocounter{table}{-1}
\normalsize
\notebox[inp]{모델에 들어간 교통량은 모든 도로가 가정값이다. 이번 조사에서 두 측정점 앞 도로(사가정로 동측, 전농로 남측)의 교통량을 측정과 동시에 계수한다(7장). 계수하지 못한 도로는 0으로 두지 않고 `미계수'로 표시한다.}
\needspace{8\baselineskip}\subsection{용도지역}
용도지역은 확인하지 못했다(미확인). 토지이음(eum.go.kr) 토지이용계획 열람으로 각 측정점 필지의 용도지역과 학교·병원 등 50 m 규정 해당 여부를 확인해 아래 표에 적는다.
\small
\begin{longtable}{>{\centering\arraybackslash}p{\dimexpr 0.1000\linewidth-2\tabcolsep\relax}>{\centering\arraybackslash}p{\dimexpr 0.2000\linewidth-2\tabcolsep\relax}>{\centering\arraybackslash}p{\dimexpr 0.2000\linewidth-2\tabcolsep\relax}>{\centering\arraybackslash}p{\dimexpr 0.1800\linewidth-2\tabcolsep\relax}>{\centering\arraybackslash}p{\dimexpr 0.1600\linewidth-2\tabcolsep\relax}>{\centering\arraybackslash}p{\dimexpr 0.1600\linewidth-2\tabcolsep\relax}}
\toprule
\hd{측정점} & \hd{필지(지번)} & \hd{용도지역} & \hd{학교·병원 50 m} & \hd{별표 1 지역} & \hd{별표 11 구분} \\ \midrule
\endfirsthead
\toprule
\hd{측정점} & \hd{필지(지번)} & \hd{용도지역} & \hd{학교·병원 50 m} & \hd{별표 1 지역} & \hd{별표 11 구분} \\ \midrule
\endhead
\bottomrule
\endlastfoot
\hr{10mm}A1 & \ip{} & \ip{} & \ip{} & \ip{} & \ip{} \\ \midrule
\hr{10mm}A2 & \ip{} & \ip{} & \ip{} & \ip{} & \ip{} \\
\end{longtable}
\addtocounter{table}{-1}
\normalsize
\needspace{8\baselineskip}\subsection{주요 수음 대상과 도로변지역}
주거(공동주택·단독주택), 학교, 병원, 어린이집 등은 현장 답사 때 확인한다. 학교·종합병원·공공도서관 부지경계 50 m 이내는 적용 기준이 ``가''지역(별표 1) 또는 가목(별표 11)으로 바뀐다. 도로변지역은 도로단에서 `차선수 $\times$ 10 m' 안이다(ES 03301.1c 5.1.1 주 1): A1 앞 사가정로 4차로 → 40 m, A2 앞 전농로 5차로 → 50 m. 두 측정점 모두 도로변지역이다(현장에서 차선수 다시 확인).
\needspace{10\baselineskip}\section{측정점 계획}
\needspace{8\baselineskip}\subsection{선정 기준}
\footnotesize
\begin{longtable}{>{\raggedright\arraybackslash}p{\dimexpr 0.1500\linewidth-2\tabcolsep\relax}>{\raggedright\arraybackslash}p{\dimexpr 0.5700\linewidth-2\tabcolsep\relax}>{\raggedright\arraybackslash}p{\dimexpr 0.2800\linewidth-2\tabcolsep\relax}}
\toprule
\hd{구분} & \hd{기준} & \hd{출처} \\ \midrule
\endfirsthead
\toprule
\hd{구분} & \hd{기준} & \hd{출처} \\ \midrule
\endhead
\bottomrule
\endlastfoot
도로 & 1시간 교통량 1,000대 이상인 도로당 한 지점 이상 → 사가정로 1, 전농로 1 & 소음지도 작성방법 별표 2 제15호 \\ \midrule
건물 & 피해가 예상되는 사람이 거주하는 건물(주거) & ES 03304.1d 5.1.1 \\ \midrule
위치 & 건물 외벽에서 도로(소음원) 방향 1.0 m. 모델 외벽 예측 위치와 같음 & ES 03301.1c 5.1.3, 별표 2 제14호 라목 \\ \midrule
높이 & 지면 위 1.2--1.5 m → 1.5 m & 두 측정방법 5.1 \\ \midrule
지점 수 & 2개 이상 & ES 03304.1d 5.3.1 \\ \midrule
자체 기준 & 간선도로에서 처음 보이는 맨 앞 건물. 수치지도 건물 용도 `주택'(아파트·주택). 기존 모델에서 외벽 소음도가 높은 쪽 우선(가정 교통량, 참고) & 본 조사 \\
\end{longtable}
\addtocounter{table}{-1}
\normalsize
\needspace{8\baselineskip}\subsection{측정점}
\footnotesize
\begin{longtable}{>{\centering\arraybackslash}p{\dimexpr 0.0600\linewidth-2\tabcolsep\relax}>{\raggedright\arraybackslash}p{\dimexpr 0.4200\linewidth-2\tabcolsep\relax}>{\centering\arraybackslash}p{\dimexpr 0.1600\linewidth-2\tabcolsep\relax}>{\centering\arraybackslash}p{\dimexpr 0.1400\linewidth-2\tabcolsep\relax}>{\centering\arraybackslash}p{\dimexpr 0.0800\linewidth-2\tabcolsep\relax}>{\centering\arraybackslash}p{\dimexpr 0.1400\linewidth-2\tabcolsep\relax}}
\caption{측정점(2026-10-21 확정, 좌표는 수치지도 판독값 --- 현장 확인 필요)}\\
\toprule
\hd{ID} & \hd{위치} & \hd{WGS84(위도, 경도)} & \hd{EPSG:5186(X, Y)} & \hd{높이} & \hd{중심선 거리} \\ \midrule
\endfirsthead
\toprule
\hd{ID} & \hd{위치} & \hd{WGS84(위도, 경도)} & \hd{EPSG:5186(X, Y)} & \hd{높이} & \hd{중심선 거리} \\ \midrule
\endhead
\bottomrule
\endlastfoot
A1 & 사가정로 동쪽 접근로 남측, 전농동하우스토리 아파트(20층, 사가정로 122) 외벽에서 도로 방향 1.0 m. 교차로 중심에서 약 55 m & 37.577467\newline 127.058450 & 205163.1\newline 553103.6 & 1.5 m & 약 16.8 m \\ \midrule
A2 & 전농로 남쪽 접근로 서측, 사가정로12길 5 주택(2층) 외벽에서 도로 방향 1.0 m. 교차로 중심에서 약 75 m & 37.576881\newline 127.057503 & 205079.5\newline 553038.5 & 1.5 m & 약 23.9 m \\
\end{longtable}
\normalsize
\begin{figure}[ht]\centering
\includegraphics[width=0.86\linewidth]{측정점_후보_위치도.png}
\caption{측정점 위치도(빨강 A1·A2 = 확정, 주황 A2$'$·초록 B = 검토한 다른 안). 수치지도 v2.0 건물 용도 기준}
\end{figure}
\needspace{8\baselineskip}\subsection{현장 실측 기입표}
지도에서 읽은 값을 현장에서 다시 잰다. 스마트폰 GPS(WGS84)로 좌표를, 거리측정기·줄자로 거리를 잰다. 측정하지 못한 값은 빈칸으로 두고, 추정값은 비고에 `추정'으로 적는다.
\footnotesize
\begin{longtable}{>{\centering\arraybackslash}p{\dimexpr 0.0600\linewidth-2\tabcolsep\relax}>{\centering\arraybackslash}p{\dimexpr 0.1200\linewidth-2\tabcolsep\relax}>{\centering\arraybackslash}p{\dimexpr 0.1200\linewidth-2\tabcolsep\relax}>{\centering\arraybackslash}p{\dimexpr 0.1100\linewidth-2\tabcolsep\relax}>{\centering\arraybackslash}p{\dimexpr 0.1000\linewidth-2\tabcolsep\relax}>{\centering\arraybackslash}p{\dimexpr 0.1000\linewidth-2\tabcolsep\relax}>{\centering\arraybackslash}p{\dimexpr 0.0900\linewidth-2\tabcolsep\relax}>{\centering\arraybackslash}p{\dimexpr 0.1500\linewidth-2\tabcolsep\relax}>{\centering\arraybackslash}p{\dimexpr 0.1500\linewidth-2\tabcolsep\relax}}
\toprule
\hd{ID} & \hd{경도} & \hd{위도} & \hd{차도 끝 거리(m)} & \hd{외벽 이격(m)} & \hd{보도 폭(m)} & \hd{차선수} & \hd{건물 용도} & \hd{사진 번호} \\ \midrule
\endfirsthead
\toprule
\hd{ID} & \hd{경도} & \hd{위도} & \hd{차도 끝 거리(m)} & \hd{외벽 이격(m)} & \hd{보도 폭(m)} & \hd{차선수} & \hd{건물 용도} & \hd{사진 번호} \\ \midrule
\endhead
\bottomrule
\endlastfoot
\hr{10mm}A1 & \ip{} & \ip{} & \ip{} & \ip{} & \ip{} & \ip{} & \ip{} & \ip{} \\ \midrule
\hr{10mm}A2 & \ip{} & \ip{} & \ip{} & \ip{} & \ip{} & \ip{} & \ip{} & \ip{} \\
\end{longtable}
\addtocounter{table}{-1}
\normalsize
\needspace{8\baselineskip}\subsection{선정 사유와 현장 확인 사항}
\footnotesize
\begin{longtable}{>{\centering\arraybackslash}p{\dimexpr 0.0600\linewidth-2\tabcolsep\relax}>{\raggedright\arraybackslash}p{\dimexpr 0.4700\linewidth-2\tabcolsep\relax}>{\raggedright\arraybackslash}p{\dimexpr 0.4700\linewidth-2\tabcolsep\relax}}
\toprule
\hd{ID} & \hd{선정 사유} & \hd{현장 확인} \\ \midrule
\endfirsthead
\toprule
\hd{ID} & \hd{선정 사유} & \hd{현장 확인} \\ \midrule
\endhead
\bottomrule
\endlastfoot
A1 & 사가정로(1시간 교통량 1,000대 이상 추정)에 직접 면한 공동주택. 피해 예상 거주 건물이며 모델 벽면 소음도가 높은 쪽 & 도로폭 27 m(4차로, 수치지도). 아파트 단지 안이면 관리사무소 동의 필요. 교차로 120 m 안(RLS-19 신호 보정 구간) \\ \midrule
A2 & 전농로(1시간 교통량 1,000대 이상 추정)에 면한 맨 앞 주거 건물. 고시의 도로당 1지점 요건 & 도로폭 25 m(5차로, 수치지도). 도로 중심선에서 약 24 m 들어간 건물 — 앞쪽 공지·주차 여부 현장 확인. 교차로 120 m 안 \\
\end{longtable}
\addtocounter{table}{-1}
\normalsize
\needspace{8\baselineskip}\subsection{검토한 다른 안}
\footnotesize
\begin{longtable}{>{\raggedright\arraybackslash}p{\dimexpr 0.1600\linewidth-2\tabcolsep\relax}>{\raggedright\arraybackslash}p{\dimexpr 0.3800\linewidth-2\tabcolsep\relax}>{\raggedright\arraybackslash}p{\dimexpr 0.4600\linewidth-2\tabcolsep\relax}}
\toprule
\hd{안} & \hd{내용} & \hd{채택하지 않은 이유} \\ \midrule
\endfirsthead
\toprule
\hd{안} & \hd{내용} & \hd{채택하지 않은 이유} \\ \midrule
\endhead
\bottomrule
\endlastfoot
B. 같은 단면 양쪽 & 사가정로 같은 단면의 A1과 맞은편 사가정로 119 앞(상·하행 비교) & 전농로가 빠져 고시 `도로당 한 지점 이상'을 못 채움. 맞은편은 근린생활시설(상가)이라 ES 03304.1d의 피해 예상 거주 건물이 아님 \\ \midrule
A2$'$. 전농로 북측 & 전농로17길 6 주택(1층) 외벽 앞, 교차로에서 약 170 m & 교차로 영향이 적은 대안. 주변 노유자 시설 50 m 여부에 따라 기준이 ``가''지역으로 바뀔 수 있음. A2가 전농로 남측 본선에 있어 A2 채택 \\
\end{longtable}
\addtocounter{table}{-1}
\normalsize
\needspace{8\baselineskip}\subsection{파일과 모델 수음점}
\begin{itemize}
\item \texttt{field\_survey/plan/measurement\_points.csv}·\texttt{.geojson} --- 측정점 목록(좌표 포함). 현장 실측 후 CSV를 고치고 \texttt{survey\_points.py --from-csv}로 GeoJSON을 다시 만든다.
\item 소프트웨어 점 수음점 입력(\texttt{receivers.points.path}, README 4.5)에 \texttt{measurement\_points.csv}를 그대로 지정하면 같은 좌표·높이에서 계산하고 \texttt{point\_levels\_for\_comparison.csv}를 출력한다. 점 수음점은 자기 벽면 반사를 포함하므로 고시 제14호 라목의 외벽 예측 조건과 같다.
\item 모델의 주간·야간 평가레벨(RLS-19)은 시간대 평균이다. 측정 구간과 비교할 때는 같은 구간의 동시 교통량(시간당 대수)으로 다시 계산한 값을 쓴다.
\end{itemize}
\needspace{10\baselineskip}\section{동시측정 계획}
\needspace{8\baselineskip}\subsection{사용 장비와 인원}
조사 인원은 5--6명이며 3명씩 2팀으로 나눈다. 팀마다 소음계 1대를 쓴다. 소음계 정보는 확인 후 아래 표에 적는다.
\footnotesize
\begin{longtable}{>{\centering\arraybackslash}p{\dimexpr 0.1000\linewidth-2\tabcolsep\relax}>{\centering\arraybackslash}p{\dimexpr 0.1700\linewidth-2\tabcolsep\relax}>{\centering\arraybackslash}p{\dimexpr 0.1700\linewidth-2\tabcolsep\relax}>{\centering\arraybackslash}p{\dimexpr 0.1000\linewidth-2\tabcolsep\relax}>{\centering\arraybackslash}p{\dimexpr 0.1800\linewidth-2\tabcolsep\relax}>{\centering\arraybackslash}p{\dimexpr 0.1300\linewidth-2\tabcolsep\relax}>{\centering\arraybackslash}p{\dimexpr 0.1500\linewidth-2\tabcolsep\relax}}
\toprule
\hd{장비 ID} & \hd{팀·측정점} & \hd{장비명/모델} & \hd{등급} & \hd{교정기(모델·표시값)} & \hd{1초 \Laeq{} 기록} & \hd{일련번호} \\ \midrule
\endfirsthead
\toprule
\hd{장비 ID} & \hd{팀·측정점} & \hd{장비명/모델} & \hd{등급} & \hd{교정기(모델·표시값)} & \hd{1초 \Laeq{} 기록} & \hd{일련번호} \\ \midrule
\endhead
\bottomrule
\endlastfoot
\hr{10mm}SLM-1 & 1팀 · A1(사가정로) & \ip{} & \ip{} & \ip{} & \ip{} & \ip{} \\ \midrule
\hr{10mm}SLM-2 & 2팀 · A2(전농로) & \ip{} & \ip{} & \ip{} & \ip{} & \ip{} \\
\end{longtable}
\addtocounter{table}{-1}
\normalsize
보조장비(팀마다): 풍속계, 온·습도계, 거리측정기, 삼각대, 스마트폰(GPS·기준 시계), 영상 촬영용 카메라 --- 보유 여부 \IN
\needspace{8\baselineskip}\subsection{기본 구조}
\begin{itemize}
\item \textbf{동시 측정}: 두 팀이 같은 회차·같은 시각에 각자의 측정점(A1, A2)에서 측정한다. 두 점은 약 100 m 떨어져 있다.
\item \textbf{회차당 30분 연속 측정}: 그 회차의 등가소음도로 하고 야장에 회차마다 1행으로 적는다(ES 03304.1d 7.1.1.1 `10분 이상', ES 03301.1c 7.1.1 `5분 이상' 충족).
\item \textbf{회차}: 낮 4회, 밤 2회. 회차 시작 간격 4시간(5.4).
\item \textbf{교정}: 회차 시작 전과 끝난 뒤 교정기로 점검·교정하고 값을 야장에 적는다(매회 교정, ES 03301.1c 3.2.2).
\item \textbf{교통량}: 각 팀이 자기 측정점 앞 도로의 양방향 교통량을 5분 단위로 센다(7.1).
\end{itemize}
\needspace{8\baselineskip}\subsection{팀 구성}
\small
\begin{longtable}{>{\raggedright\arraybackslash}p{\dimexpr 0.2000\linewidth-2\tabcolsep\relax}>{\raggedright\arraybackslash}p{\dimexpr 0.8000\linewidth-2\tabcolsep\relax}}
\toprule
\hd{역할(팀마다)} & \hd{업무} \\ \midrule
\endfirsthead
\toprule
\hd{역할(팀마다)} & \hd{업무} \\ \midrule
\endhead
\bottomrule
\endlastfoot
측정자 & 소음계 설치·교정·측정, 소음 기록 카드 작성, 이상 소음 시각 기록 \\ \midrule
교통량 계수 & 측정점 앞 도로 양방향 $\times$ 6개 차종, 5분 단위(영상 녹화 병행 권장) \\ \midrule
기록·안전 & 기상(풍속·풍향·기온·습도·노면), 신호 주기·주행속도, 사진, 보행자 안내·안전 확인 \\
\end{longtable}
\addtocounter{table}{-1}
\normalsize
{\small 6명이면 팀마다 3명이다. 5명이면 한 팀은 2명이 되므로 그 팀은 교통량을 영상으로만 기록하고 현장 계수를 생략한다.\par}
\needspace{8\baselineskip}\subsection{측정 회차 시각}
회차 시작을 4시간 간격으로 두면 환경기준(낮 2시간 이상 간격 4회, 밤 2시간 간격 2회)과 도로교통 관리기준(4시간 이상 간격 2회, 야간은 00--04시 포함)을 모든 회차가 함께 만족한다. 출근 시간(D1)과 퇴근 이후(D4)가 피해 예상 시간대를 포함한다.
\small
\begin{longtable}{>{\centering\arraybackslash}p{\dimexpr 0.0900\linewidth-2\tabcolsep\relax}>{\centering\arraybackslash}p{\dimexpr 0.1400\linewidth-2\tabcolsep\relax}>{\centering\arraybackslash}p{\dimexpr 0.1200\linewidth-2\tabcolsep\relax}>{\centering\arraybackslash}p{\dimexpr 0.1200\linewidth-2\tabcolsep\relax}>{\centering\arraybackslash}p{\dimexpr 0.1300\linewidth-2\tabcolsep\relax}>{\centering\arraybackslash}p{\dimexpr 0.1400\linewidth-2\tabcolsep\relax}>{\centering\arraybackslash}p{\dimexpr 0.2600\linewidth-2\tabcolsep\relax}}
\toprule
\hd{회차} & \hd{시간대} & \hd{측정 시작} & \hd{측정 끝} & \hd{교정(후) 끝} & \hd{환경기준} & \hd{도로교통 관리기준} \\ \midrule
\endfirsthead
\toprule
\hd{회차} & \hd{시간대} & \hd{측정 시작} & \hd{측정 끝} & \hd{교정(후) 끝} & \hd{환경기준} & \hd{도로교통 관리기준} \\ \midrule
\endhead
\bottomrule
\endlastfoot
D1 & 낮(주간) & 07:30 & 08:00 & 08:10 & 사용 & 사용 \\ \midrule
D2 & 낮(주간) & 11:30 & 12:00 & 12:10 & 사용 & 사용 \\ \midrule
D3 & 낮(주간) & 15:30 & 16:00 & 16:10 & 사용 & 사용 \\ \midrule
D4 & 낮(주간) & 19:30 & 20:00 & 20:10 & 사용 & 사용 \\ \midrule
N1 & 밤(야간) & 22:30 & 23:00 & 23:10 & 사용 & 사용 \\ \midrule
N2 & 밤(야간) & 02:30 & 03:00 & 03:10 & 사용 & 사용 (00--04시 포함) \\
\end{longtable}
\addtocounter{table}{-1}
\normalsize
\notebox{야간 측정을 하지 않으면 밤 시간대 기준과의 비교(목적 가)는 `평가 불가'로 보고하고, 측정하지 않은 값을 모델값이나 추정값으로 채우지 않는다. 낮 회차를 하루에 모두 하지 못하면 회차 간 4시간 간격을 지켜 다른 평일에 이어서 하고, 날짜를 야장에 적는다.}
\needspace{8\baselineskip}\subsection{회차 내 시각표}
시각은 회차 시작 기준(+시:분)이다. 색칸은 측정 시간이며 두 팀이 같은 시각에 측정한다.
\footnotesize
\begin{longtable}{>{\centering\arraybackslash}p{\dimexpr 0.1800\linewidth-2\tabcolsep\relax}>{\centering\arraybackslash}p{\dimexpr 0.2000\linewidth-2\tabcolsep\relax}>{\centering\arraybackslash}p{\dimexpr 0.2000\linewidth-2\tabcolsep\relax}>{\centering\arraybackslash}p{\dimexpr 0.2000\linewidth-2\tabcolsep\relax}>{\centering\arraybackslash}p{\dimexpr 0.2200\linewidth-2\tabcolsep\relax}}
\caption{회차 내 시각표(두 팀 공통)}\\
\toprule
\hd{회차 시작 기준} & \hd{1팀(A1, SLM-1)} & \hd{2팀(A2, SLM-2)} & \hd{교통량(각 팀)} & \hd{기상·기록(각 팀)} \\ \midrule
\endfirsthead
\toprule
\hd{회차 시작 기준} & \hd{1팀(A1, SLM-1)} & \hd{2팀(A2, SLM-2)} & \hd{교통량(각 팀)} & \hd{기상·기록(각 팀)} \\ \midrule
\endhead
\bottomrule
\endlastfoot
$-$0:15 -- $-$0:05 & 교정(전)·설치 & 교정(전)·설치 & 시각 맞춤·위치 잡기 & 기상 1회 \\ \midrule
+0:00 -- +0:10 & \cellcolor{team1}{}\textbf{측정 ①} & \cellcolor{team2}{}\textbf{측정 ①} & 5분 계수 $\times$2 & 이상 소음·신호주기 \\ \midrule
+0:10 -- +0:20 & \cellcolor{team1}{}\textbf{측정 ②} & \cellcolor{team2}{}\textbf{측정 ②} & 5분 계수 $\times$2 & 이상 소음·주행속도 \\ \midrule
+0:20 -- +0:30 & \cellcolor{team1}{}\textbf{측정 ③} & \cellcolor{team2}{}\textbf{측정 ③} & 5분 계수 $\times$2 & 이상 소음·기상 1회 \\ \midrule
+0:30 -- +0:40 & 교정(후)·저장 & 교정(후)·저장 & 계수표 정리 & 야장 확인 \\
\end{longtable}
\normalsize
\needspace{8\baselineskip}\subsection{시각 동기화}
\begin{itemize}
\item 측정일 아침, 소음계·카메라·스마트폰의 시계를 같은 기준 시계(네트워크 시각 스마트폰)에 초 단위로 맞춘다.
\item 맞춘 직후와 하루 끝에 각 장비 화면과 기준 시계를 함께 사진으로 찍고, 차이(초)를 야장 `시각 오차'에 적는다.
\item 카메라는 녹화 시작 직후 기준 시계 화면을 10초 동안 찍는다. 두 팀은 회차 시작 15분 전에 서로 연락해 시작 시각을 맞춘다.
\end{itemize}
\needspace{10\baselineskip}\section{측정 조건}
\needspace{8\baselineskip}\subsection{측정기 설정}
\footnotesize
\begin{longtable}{>{\raggedright\arraybackslash}p{\dimexpr 0.1800\linewidth-2\tabcolsep\relax}>{\raggedright\arraybackslash}p{\dimexpr 0.4200\linewidth-2\tabcolsep\relax}>{\raggedright\arraybackslash}p{\dimexpr 0.4000\linewidth-2\tabcolsep\relax}}
\toprule
\hd{항목} & \hd{설정} & \hd{근거} \\ \midrule
\endfirsthead
\toprule
\hd{항목} & \hd{설정} & \hd{근거} \\ \midrule
\endhead
\bottomrule
\endlastfoot
청감보정 & A 특성 & ES 03301.1c 3.3.1, ES 03304.1d 3.3.1 \\ \midrule
동특성 & 빠름(fast) & ES 03301.1c 3.3.2, ES 03304.1d 3.3.2 \\ \midrule
샘플주기 & 1초 이하 & ES 03301.1c 7.1.1, ES 03304.1d 3.3.3 \\ \midrule
측정시간 & 회차당 30분 연속 & ES 03301.1c 7.1.1(5분 이상), ES 03304.1d 7.1.1.1(10분 이상) \\ \midrule
레벨레인지 & 예비조사 후 고정 & ES 03301.1c 3.2.3 \\ \midrule
저장 지표 & \Laeq, $L_{\mathrm{Amax}}$, $L_{\mathrm{Amin}}$, $L_{10}$, $L_{50}$, $L_{90}$ 중 장비가 제공하는 값 & 자체(기록 항목) \\ \midrule
마이크로폰 & 삼각대, 주 소음원(도로) 방향, 방풍망 부착 & 두 방법 5.2.1 \\
\end{longtable}
\addtocounter{table}{-1}
\normalsize
\needspace{8\baselineskip}\subsection{교정 절차}
\begin{itemize}
\item 회차 시작 전과 끝난 뒤 교정기를 씌워 지시값을 읽고 교정한다. 교정 전 지시값과 교정 후 값을 함께 적는다.
\item 교정기는 발생음 주파수·음압도가 표시되고 오차 $\pm$1 dB 이내인 것을 쓴다(ES 03300.c 3.2.3).
\item 교정 전후 허용 편차는 적용 방법에 규정이 없다(미확인). 본 조사의 판정 기준은 \IN\ dB로 하며 지도교수 확인 후 정한다. 넘으면 그 회차 값은 `무효 후보'로 표시한다.
\end{itemize}
\needspace{8\baselineskip}\subsection{기상 조건}
\footnotesize
\begin{longtable}{>{\raggedright\arraybackslash}p{\dimexpr 0.1800\linewidth-2\tabcolsep\relax}>{\raggedright\arraybackslash}p{\dimexpr 0.6000\linewidth-2\tabcolsep\relax}>{\raggedright\arraybackslash}p{\dimexpr 0.2200\linewidth-2\tabcolsep\relax}}
\toprule
\hd{조건} & \hd{조치} & \hd{근거} \\ \midrule
\endfirsthead
\toprule
\hd{조건} & \hd{조치} & \hd{근거} \\ \midrule
\endhead
\bottomrule
\endlastfoot
풍속 2 m/s 이상 & 방풍망 필수(본 조사는 항상 부착). 야장에 방풍망 Y/N 기록 & 두 방법 5.2.1.4 \\ \midrule
풍속 5 m/s 초과 & 측정 중단. 그 구간은 무효, 다른 날 같은 시각에 다시 측정 & 두 방법 5.2.1.4 \\ \midrule
강우·노면 젖음 & 원문 규정 없음(미확인). 비가 오거나 노면이 젖으면 측정을 미루고 노면 상태를 기록(자체) & 자체 \\ \midrule
기록 & 회차마다 풍속(마이크로폰 높이)·풍향·기온·습도·노면 상태 & 자체 \\
\end{longtable}
\addtocounter{table}{-1}
\normalsize
\needspace{8\baselineskip}\subsection{배경소음}
\begin{itemize}
\item 법정 평가: 환경기준·도로교통 방법에는 배경소음 보정 조항이 없으므로 보정하지 않는다.
\item 검증용: 도로교통 이외의 이상소음(경적, 사이렌, 공사 등)이 들어간 시각을 기록하고, 1초 자료가 있으면 그 구간을 뺀 값도 계산해 두 값을 모두 보고한다(자체).
\end{itemize}
\needspace{10\baselineskip}\section{동시 기록 항목(시뮬레이션 검증용)}
\needspace{8\baselineskip}\subsection{교통량}
\begin{itemize}
\item 계수 단위: 5분. 방향별 $\times$ 차종별. 이륜차는 따로 센다.
\item 계수 대상: 1팀은 A1 앞 사가정로 동측 접근로 양방향, 2팀은 A2 앞 전농로 남측 접근로 양방향. 그 밖의 접근로는 인원이 없어 세지 않으며 모델에서 가정값으로 표시한다.
\item 방법: 측정점 근처에서 도로가 보이는 곳에 카메라를 두고 녹화하면서 현장에서도 센다. 영상은 계수 확인용으로만 쓴다.
\item 계수하지 못한 도로·시간은 `미계수'로 표시하고 0으로 적지 않는다.
\end{itemize}
\small
\begin{longtable}{>{\raggedright\arraybackslash}p{\dimexpr 0.2200\linewidth-2\tabcolsep\relax}>{\raggedright\arraybackslash}p{\dimexpr 0.4800\linewidth-2\tabcolsep\relax}>{\raggedright\arraybackslash}p{\dimexpr 0.3000\linewidth-2\tabcolsep\relax}}
\caption{계수 차종과 소프트웨어 입력 차종(RLS-19 차량군) 대응}\\
\toprule
\hd{계수 차종} & \hd{예} & \hd{RLS-19 차량군} \\ \midrule
\endfirsthead
\toprule
\hd{계수 차종} & \hd{예} & \hd{RLS-19 차량군} \\ \midrule
\endhead
\bottomrule
\endlastfoot
1 승용·승합 & 승용차, SUV, 15인 미만 승합 & Pkw \\ \midrule
2 소·중형 화물 & 3.5 t 이하 화물(1톤·2.5톤 트럭, 밴) & Pkw \\ \midrule
3 버스 & 15인 이상 버스(시내·마을·광역·관광) & Lkw1 \\ \midrule
4 대형 화물 & 3.5 t 초과, 트레일러 없음(카고·덤프·탱크) & Lkw1 \\ \midrule
5 트레일러 & 트레일러·세미트레일러 연결차 & Lkw2 \\ \midrule
6 이륜차 & 오토바이·스쿠터(배달 포함) & Lkw2(방출만, 속도는 Pkw) \\
\end{longtable}
\normalsize
{\small 근거: RLS-19 2장 차량군 정의(Lkw1 = 허용 총중량 3.5 t 초과 트레일러 없는 화물차와 버스, Lkw2 = 트레일러 부착 화물차·세미트레일러 연결차, 이륜자동차는 방출 측면에서 Lkw2로 분류하고 속도는 Pkw 속도)와 「소음지도의 작성방법」 별표 2 제12호 다목. 현장에서는 3.5 t 경계를 차체 크기로 판단하고 애매한 차량은 메모에 적는다.\par}
\needspace{8\baselineskip}\subsection{신호 주기·정체·주행속도}
\begin{itemize}
\item 신호 주기: 회차마다 2주기 이상, 스톱워치로 주기 길이와 녹색 시간을 잰다.
\item 정체: 5분마다 `원활 / 서행 / 정체'를 적는다.
\item 주행속도(자체 방법): 보도에 거리 \IN\ m 구간을 정해 신호에 막히지 않은 차량 10대 이상의 통과 시간을 재어 $v = 3.6 \times$ 거리(m) / 시간(s) [km/h]로 계산한다.
\end{itemize}
\needspace{8\baselineskip}\subsection{이상 소음 기록}
\small
\begin{longtable}{>{\centering\arraybackslash}p{\dimexpr 0.1000\linewidth-2\tabcolsep\relax}>{\raggedright\arraybackslash}p{\dimexpr 0.4000\linewidth-2\tabcolsep\relax}>{\centering\arraybackslash}p{\dimexpr 0.1000\linewidth-2\tabcolsep\relax}>{\raggedright\arraybackslash}p{\dimexpr 0.4000\linewidth-2\tabcolsep\relax}}
\toprule
\hd{코드} & \hd{내용} & \hd{코드} & \hd{내용} \\ \midrule
\endfirsthead
\toprule
\hd{코드} & \hd{내용} & \hd{코드} & \hd{내용} \\ \midrule
\endhead
\bottomrule
\endlastfoot
H & 경적 & S & 사이렌(구급·경찰·소방) \\ \midrule
M & 이륜차 급가속·개조 배기음 & C & 공사·작업 소음 \\ \midrule
V & 사람 대화·확성기·음악 & A & 항공기 \\ \midrule
B & 버스 정차·문 열림 소리 & O & 기타(메모) \\
\end{longtable}
\addtocounter{table}{-1}
\normalsize
\needspace{8\baselineskip}\subsection{사진·영상 기록}
\begin{itemize}
\item 측정점마다: 마이크로폰 설치 모습 1장, 측정점에서 동·서·남·북 4장, 외벽과 거리 표시 1장.
\item 사람 얼굴·차량 번호판이 알아볼 수 있게 찍히지 않도록 한다. 교통 영상은 계수에만 쓰고 공개하지 않는다.
\item 아파트 단지·사유지 안은 관리사무소·소유자 동의를 받은 뒤 들어가고 찍는다.
\end{itemize}
\needspace{10\baselineskip}\section{안전 및 현장 수칙}
\begin{itemize}
\item 모든 장비와 사람은 보도 위에 둔다. 차도·중앙분리대에 들어가지 않는다.
\item 반사 조끼를 입는다. 야간에는 헤드랜턴과 삼각대 반사 띠를 쓰고 2인 1조로 움직인다.
\item 삼각대와 케이블이 보행을 막지 않게 두고 고정한다. 장비 옆을 비우지 않는다.
\item 사유지는 동의를 받은 뒤 들어간다. 주민이 물으면 조사 목적(학교 과제)을 설명하고 개인 정보는 기록하지 않는다.
\item 기상 악화 시 즉시 중단하고 장비를 거둔다.
\end{itemize}
\needspace{10\baselineskip}\section{역할 분담, 준비물, 대체 일정}
\needspace{8\baselineskip}\subsection{역할 분담}
\small
\begin{longtable}{>{\centering\arraybackslash}p{\dimexpr 0.1000\linewidth-2\tabcolsep\relax}>{\centering\arraybackslash}p{\dimexpr 0.2700\linewidth-2\tabcolsep\relax}>{\centering\arraybackslash}p{\dimexpr 0.2100\linewidth-2\tabcolsep\relax}>{\centering\arraybackslash}p{\dimexpr 0.2100\linewidth-2\tabcolsep\relax}>{\centering\arraybackslash}p{\dimexpr 0.2100\linewidth-2\tabcolsep\relax}}
\toprule
\hd{팀} & \hd{측정점} & \hd{측정자} & \hd{교통량 계수} & \hd{기록·안전} \\ \midrule
\endfirsthead
\toprule
\hd{팀} & \hd{측정점} & \hd{측정자} & \hd{교통량 계수} & \hd{기록·안전} \\ \midrule
\endhead
\bottomrule
\endlastfoot
\hr{10mm}1팀 & A1 사가정로 (SLM-1) & \ip{} & \ip{} & \ip{} \\ \midrule
\hr{10mm}2팀 & A2 전농로 (SLM-2) & \ip{} & \ip{} & \ip{} \\
\end{longtable}
\addtocounter{table}{-1}
\normalsize
{\small 총괄(시각 동기화, 회차 진행 연락)은 1팀 기록 담당이 맡는다.\par}
\needspace{8\baselineskip}\subsection{준비물 체크리스트}
\footnotesize
\begin{longtable}{>{\centering\arraybackslash}p{\dimexpr 0.0700\linewidth-2\tabcolsep\relax}>{\raggedright\arraybackslash}p{\dimexpr 0.2600\linewidth-2\tabcolsep\relax}>{\centering\arraybackslash}p{\dimexpr 0.0700\linewidth-2\tabcolsep\relax}>{\raggedright\arraybackslash}p{\dimexpr 0.2600\linewidth-2\tabcolsep\relax}>{\centering\arraybackslash}p{\dimexpr 0.0700\linewidth-2\tabcolsep\relax}>{\raggedright\arraybackslash}p{\dimexpr 0.2700\linewidth-2\tabcolsep\relax}}
\toprule
\hd{확인} & \hd{준비물} & \hd{확인} & \hd{준비물} & \hd{확인} & \hd{준비물} \\ \midrule
\endfirsthead
\toprule
\hd{확인} & \hd{준비물} & \hd{확인} & \hd{준비물} & \hd{확인} & \hd{준비물} \\ \midrule
\endhead
\bottomrule
\endlastfoot
$\square$ & 소음계(예비 배터리) & $\square$ & 교정기 & $\square$ & 방풍망 \\ \midrule
$\square$ & 삼각대 & $\square$ & 풍속계 & $\square$ & 온·습도계 \\ \midrule
$\square$ & 거리측정기·줄자 & $\square$ & 스마트폰(GPS·기준 시계) & $\square$ & 카메라·카메라 삼각대 \\ \midrule
$\square$ & 보조배터리 & $\square$ & 스톱워치 & $\square$ & 인쇄한 야장·클립보드·필기구 \\ \midrule
$\square$ & 측정점 지도 출력본 & $\square$ & 반사 조끼 & $\square$ & 헤드랜턴(야간) \\ \midrule
$\square$ & 우비·장비 덮개 & $\square$ & 모래주머니·테이프·분필 & $\square$ & 동의서 양식 \\
\end{longtable}
\addtocounter{table}{-1}
\normalsize
\needspace{8\baselineskip}\subsection{우천·강풍 시 대체 일정}
\begin{itemize}
\item 측정 전날 \IN 시 일기예보로 진행 여부를 정한다. 비 예보 또는 풍속 5 m/s 초과 예보면 예비일 \IN(평일)로 미룬다.
\item 측정 중 비가 오거나 풍속이 5 m/s를 넘으면 그 회차를 중단하고, 남은 회차는 4시간 간격을 지켜 예비일 같은 시각에 이어서 한다.
\end{itemize}
\needspace{10\baselineskip}\section{데이터 관리}
\footnotesize
\begin{longtable}{>{\raggedright\arraybackslash}p{\dimexpr 0.2000\linewidth-2\tabcolsep\relax}>{\raggedright\arraybackslash}p{\dimexpr 0.4200\linewidth-2\tabcolsep\relax}>{\raggedright\arraybackslash}p{\dimexpr 0.3800\linewidth-2\tabcolsep\relax}}
\toprule
\hd{종류} & \hd{파일명 규칙} & \hd{예} \\ \midrule
\endfirsthead
\toprule
\hd{종류} & \hd{파일명 규칙} & \hd{예} \\ \midrule
\endhead
\bottomrule
\endlastfoot
소음계 원자료 & JN\_날짜\_회차\_측정점\_장비.확장자 & JN\_20261021\_D1\_A1\_SLM1.csv \\ \midrule
사진 & JN\_날짜\_측정점\_방향\_번호.jpg & JN\_20261021\_A1\_N\_01.jpg \\ \midrule
교통 영상 & JN\_날짜\_회차\_카메라.mp4 & JN\_20261021\_D1\_CAM1.mp4 \\ \midrule
야장 & 소음측정\_야장\_날짜.xlsx / 야장스캔\_날짜\_쪽.pdf & 소음측정\_야장\_20261021.xlsx \\
\end{longtable}
\addtocounter{table}{-1}
\normalsize
\begin{itemize}
\item 원자료는 \texttt{field\_survey/data/raw/날짜/}에 두고 고치지 않는다. 처리 결과는 \texttt{process\_measurements.py}로만 만든다.
\item 모든 값에 근거(basis)를 적는다: measured = 직접 측정·계수, estimated = 측정값에서 환산·지도 판독, assumed = 측정 없이 가정.
\end{itemize}
\clearpage\section*{부록 A. 사용자가 채워야 할 항목}\addcontentsline{toc}{section}{부록 A. 사용자가 채워야 할 항목}
\small
\begin{longtable}{>{\raggedright\arraybackslash}p{\dimexpr 0.6600\linewidth-2\tabcolsep\relax}>{\raggedright\arraybackslash}p{\dimexpr 0.3400\linewidth-2\tabcolsep\relax}}
\toprule
\hd{항목} & \hd{위치} \\ \midrule
\endfirsthead
\toprule
\hd{항목} & \hd{위치} \\ \midrule
\endhead
\bottomrule
\endlastfoot
측정 예정일·예비일, 야간 측정 여부 & 표지, 1.2, 5.4, 9.3 \\ \midrule
소음계·보조장비 현황 & 5.1 \\ \midrule
측정점 현장 실측(좌표·거리·용도), 사유지 동의 & 4.3, 4.4, measurement\_points.csv \\ \midrule
용도지역, 학교·병원 50 m 해당 여부 & 3.3 \\ \midrule
교정 전후 허용 편차 & 6.2 \\ \midrule
주행속도 측정 구간 거리 & 7.2 \\ \midrule
역할 분담 담당자 & 9.1 \\
\end{longtable}
\addtocounter{table}{-1}
\normalsize

\end{document}
